% Compilar desde este directorio con:
% pdflatex Clase_01_Fundamentos_de_IA_agentes_y_racionalidad.tex
\documentclass[aspectratio=169,11pt]{beamer}

\usepackage[utf8]{inputenc}
\usepackage[T1]{fontenc}
\usepackage[spanish,es-nodecimaldot]{babel}
\usepackage{amsmath}
\usepackage{amssymb}
\usepackage{booktabs}
\usepackage{graphicx}
\usepackage{tabularx}
\usepackage{array}

\definecolor{UAPNavy}{HTML}{17324D}
\definecolor{UAPTeal}{HTML}{007F82}
\definecolor{UAPGold}{HTML}{E2A33A}
\definecolor{UAPLight}{HTML}{F3F6F8}
\definecolor{UAPGray}{HTML}{536471}

\usetheme{Boadilla}
\usecolortheme{default}
\setbeamercolor{normal text}{fg=UAPNavy,bg=white}
\setbeamercolor{structure}{fg=UAPTeal}
\setbeamercolor{frametitle}{fg=white,bg=UAPNavy}
\setbeamercolor{title}{fg=white}
\setbeamercolor{subtitle}{fg=UAPGold}
\setbeamercolor{block title}{fg=white,bg=UAPTeal}
\setbeamercolor{block body}{fg=UAPNavy,bg=UAPLight}
\setbeamercolor{alerted text}{fg=UAPTeal}
\setbeamerfont{frametitle}{series=\bfseries}
\setbeamertemplate{navigation symbols}{}
\setbeamertemplate{itemize item}{\color{UAPGold}$\blacktriangleright$}
\setbeamertemplate{itemize subitem}{\color{UAPTeal}$\bullet$}
\setbeamertemplate{footline}{%
  \leavevmode%
  \hbox{%
    \begin{beamercolorbox}[wd=.78\paperwidth,ht=2.8ex,dp=1.2ex,leftskip=1em]{author in head/foot}%
      \color{UAPGray}\insertshorttitle
    \end{beamercolorbox}%
    \begin{beamercolorbox}[wd=.22\paperwidth,ht=2.8ex,dp=1.2ex,rightskip=1em plus1fil]{date in head/foot}%
      \color{UAPGray}\hfill Semana 1 \enspace | \enspace \insertframenumber/\inserttotalframenumber
    \end{beamercolorbox}%
  }%
}

\newcommand{\concepto}[1]{\textcolor{UAPTeal}{\textbf{#1}}}
\newcommand{\pregunta}[1]{\begin{block}{Pregunta clave}#1\end{block}}

\title[IA · Clase 1]{Fundamentos de Inteligencia Artificial}
\subtitle{Enfoques, agentes y racionalidad}
\author{Curso de Inteligencia Artificial}
\institute{Semana 1 · Clase 1}
\date{}

\begin{document}

\begin{frame}[plain]
  \begin{beamercolorbox}[wd=\paperwidth,ht=\paperheight,sep=2em]{frametitle}
    \vfill
    {\usebeamerfont{title}\usebeamercolor[fg]{title}\Huge\inserttitle\par}
    \vspace{0.5em}
    {\usebeamerfont{subtitle}\usebeamercolor[fg]{subtitle}\Large\insertsubtitle\par}
    \vspace{2em}
    {\color{white}\large\insertauthor\par}
    {\color{white}\insertinstitute\par}
    \vfill
  \end{beamercolorbox}
\end{frame}

\begin{frame}{Propósito de la clase}
  Comprender distintas formas de definir la Inteligencia Artificial y utilizarlas para especificar un agente situado en un entorno.
  \vspace{0.5em}
  \begin{block}{Al finalizar, podremos}
    \begin{itemize}
      \item distinguir los cuatro enfoques clásicos de la IA;
      \item reconocer hitos, continuidades y límites de su evolución;
      \item explicar agente, percepción, acción y racionalidad;
      \item especificar una tarea mediante PEAS;
      \item caracterizar el ambiente de un agente de movilidad urbana.
    \end{itemize}
  \end{block}
\end{frame}

\begin{frame}{Pregunta de apertura}
  \pregunta{¿Cuándo podemos afirmar que un sistema se comporta de manera inteligente?}
  \begin{itemize}
    \item ¿Imita correctamente a una persona?
    \item ¿Utiliza procesos similares al pensamiento humano?
    \item ¿Deriva conclusiones válidas?
    \item ¿Elige acciones adecuadas para sus objetivos?
  \end{itemize}
  \vfill
  \centering\alert{La respuesta depende del criterio de inteligencia que adoptemos.}
\end{frame}

\begin{frame}{¿Qué estudia la Inteligencia Artificial?}
  La IA estudia la construcción de sistemas capaces de:
  \begin{center}
    \concepto{percibir · representar · razonar · aprender · planificar · comunicar · actuar}
  \end{center}
  En sentido técnico, ``inteligente'' significa transformar información en comportamiento dirigido a un objetivo bajo ciertas condiciones.
  \vspace{0.5em}
  \begin{block}{No implica necesariamente}
    Conciencia, comprensión humana general, intenciones propias ni autonomía completa.
  \end{block}
\end{frame}

\begin{frame}{IA, aprendizaje automático y Ciencia de Datos}
  \begin{itemize}
    \item La \concepto{IA} construye sistemas que perciben, razonan, aprenden o actúan.
    \item El \concepto{aprendizaje automático} ajusta funciones o políticas desde datos o experiencia.
    \item La \concepto{Ciencia de Datos} produce evidencia para apoyar decisiones.
  \end{itemize}
  \vspace{0.5em}
  \begin{block}{Relaciones importantes}
    \begin{itemize}
      \item Puede existir IA basada en reglas sin aprendizaje automático.
      \item Puede existir análisis predictivo sin un agente autónomo.
      \item Un sistema puede combinar reglas, búsqueda, probabilidad y aprendizaje.
    \end{itemize}
  \end{block}
\end{frame}

\begin{frame}{Cuatro enfoques clásicos}
  \centering
  \renewcommand{\arraystretch}{1.5}
  \begin{tabular}{>{\bfseries}c|c|c}
    & Como humanos & Racionalmente \\
    \hline
    Pensar & Modelado cognitivo & Leyes del pensamiento \\
    Actuar & Test de Turing & Agente racional \\
  \end{tabular}
  \vfill
  \begin{itemize}
    \item ¿Qué conducta observamos?
    \item ¿Qué proceso produce la respuesta?
    \item ¿La inferencia es válida?
    \item ¿La acción mejora el desempeño esperado?
  \end{itemize}
\end{frame}

\begin{frame}{Actuar como humanos}
  Evalúa la \concepto{conducta observable} del sistema en una tarea definida.
  \begin{columns}[T]
    \column{0.48\textwidth}
      \begin{block}{Puede considerar}
        \begin{itemize}
          \item conversación y lenguaje;
          \item percepción;
          \item movimiento;
          \item interacción.
        \end{itemize}
      \end{block}
    \column{0.48\textwidth}
      \begin{block}{Alcance}
        Permite evaluar resultados sin definir el mecanismo interno.

        \vspace{0.5em}
        Imitar una respuesta no demuestra comprensión ni confiabilidad.
      \end{block}
  \end{columns}
\end{frame}

\begin{frame}{El test de Turing}
  Una persona conversa mediante un canal textual con interlocutores ocultos e intenta distinguir a la máquina de la persona.
  \begin{block}{Aporte histórico}
    Desplaza la pregunta ``¿puede pensar una máquina?'' hacia una prueba de comportamiento.
  \end{block}
  \begin{block}{Límites}
    \begin{itemize}
      \item Evalúa una interacción particular, no inteligencia general.
      \item No garantiza verdad, seguridad ni comprensión causal.
      \item Puede favorecer imitación, evasión o plausibilidad.
      \item No evalúa necesariamente percepción física o planificación prolongada.
    \end{itemize}
  \end{block}
\end{frame}

\begin{frame}{Pensar como humanos}
  Busca modelar procesos de la cognición:
  \begin{center}
    \concepto{percepción · memoria · aprendizaje · lenguaje · razonamiento}
  \end{center}
  La evaluación observa no solo la respuesta final, sino también:
  \begin{itemize}
    \item tiempos de respuesta y secuencia de pasos;
    \item patrones de error;
    \item reacción ante información nueva.
  \end{itemize}
  \vfill
  \centering\alert{Resolver correctamente una tarea no demuestra que se utilizó un proceso humano.}
\end{frame}

\begin{frame}{Pensar racionalmente}
  Consiste en derivar conclusiones válidas desde premisas y reglas.
  \begin{block}{Ejemplo de inferencia}
    Si hay congestión, aumenta el tiempo de viaje.\\
    Hay congestión.\\
    $\therefore$ Aumenta el tiempo de viaje.
  \end{block}
  \begin{columns}[T]
    \column{0.48\textwidth}
      \textbf{Fortaleza:} premisas, reglas y conclusiones pueden inspeccionarse.
    \column{0.48\textwidth}
      \textbf{Límite:} el mundo real contiene ruido, ambigüedad y excepciones.
  \end{columns}
\end{frame}

\begin{frame}{Razonar bajo incertidumbre}
  La lógica clásica trabaja con proposiciones verdaderas o falsas. La probabilidad representa grados de creencia:
  \[
    P(\text{estado}\mid\text{evidencia})
  \]
  \begin{block}{Ejemplo}
    \centering $P(\text{congestión}\mid\text{velocidad, hora, zona}) = 0{,}75$
  \end{block}
  La evidencia modifica nuestras creencias, pero aún falta decidir:
  \begin{itemize}
    \item ¿Qué acciones están disponibles?
    \item ¿Qué consecuencias tiene cada acción?
    \item ¿Qué errores son más costosos?
  \end{itemize}
\end{frame}

\begin{frame}{Actuar racionalmente}
  Un sistema actúa racionalmente cuando elige la acción con mejor desempeño esperado según sus percepciones, objetivos y limitaciones.
  \[
    a^* = \underset{a}{\operatorname{arg\,max}}\;
    \mathbb{E}\!\left[U(a,S)\mid\text{evidencia}\right]
  \]
  \begin{itemize}
    \item $a$: acción disponible.
    \item $S$: estado posible del entorno.
    \item $U$: utilidad o medida de desempeño.
  \end{itemize}
  \begin{block}{Idea central}
    La acción racional depende de la información disponible, no de conocer el futuro con certeza.
  \end{block}
\end{frame}

\begin{frame}{Racionalidad no es perfección}
  Un agente racional:
  \begin{itemize}
    \item puede equivocarse bajo incertidumbre;
    \item no conoce aquello que no puede observar;
    \item dispone de recursos limitados;
    \item puede obtener información antes de actuar;
    \item evalúa resultados esperados, no garantías absolutas.
  \end{itemize}
  \vfill
  \begin{alertblock}{Advertencia}
    Un agente puede ejecutar perfectamente una función mal definida. Optimizar solo rapidez puede producir respuestas rápidas pero incorrectas.
  \end{alertblock}
\end{frame}

\begin{frame}{Fundamentos de la IA}
  \scriptsize
  \begin{tabularx}{\textwidth}{>{\bfseries}p{0.25\textwidth} X}
    \toprule
    Fundamento & Aporte \\
    \midrule
    Filosofía & Razón, conocimiento, experiencia y consecuencias \\
    Lógica & Representación explícita e inferencia \\
    Probabilidad & Incertidumbre y actualización de creencias \\
    Teoría de decisión & Utilidad, costos y elección de acciones \\
    Optimización & Selección de parámetros, planes o políticas \\
    Matemática & Álgebra lineal, cálculo e información \\
    Computación & Algoritmos, memoria, complejidad y sistemas \\
    \bottomrule
  \end{tabularx}
\end{frame}

\begin{frame}{Anatomía de un sistema inteligente}
  \begin{enumerate}
    \item percepción o entrada de datos;
    \item representación interna;
    \item memoria o estado;
    \item inferencia, búsqueda o aprendizaje;
    \item evaluación de acciones;
    \item ejecución y retroalimentación.
  \end{enumerate}
  \vfill
  \pregunta{¿Falló la percepción, la representación, la predicción o la decisión?}
\end{frame}

\begin{frame}{Evolución: ideas, recursos y expectativas}
  \scriptsize
  \begin{tabularx}{\textwidth}{>{\bfseries}p{0.20\textwidth} X X}
    \toprule
    Periodo & Enfoque destacado & Limitación visible \\
    \midrule
    Antecedentes & Lógica, probabilidad y algoritmos & Cálculo y representación \\
    Década de 1950 & Búsqueda y símbolos & Explosión de estados \\
    Sistemas expertos & Reglas y conocimiento & Adquisición y mantenimiento \\
    Resurgimiento estadístico & Aprendizaje desde datos & Sesgo y generalización \\
    Aprendizaje profundo & Representaciones aprendidas & Datos, cómputo y explicación \\
    IA generativa & Modelos fundacionales & Veracidad, control y uso \\
    \bottomrule
  \end{tabularx}
\end{frame}

\begin{frame}{Nacimiento formal y enfoque simbólico}
  \begin{itemize}
    \item Dartmouth, 1956, es un hito nominal del campo.
    \item Los primeros proyectos abordaron teoremas, juegos y problemas acotados.
    \item Los sistemas expertos separaron conocimiento y motor de inferencia.
    \item Las representaciones simbólicas permiten inspeccionar reglas y conclusiones.
  \end{itemize}
  \begin{block}{Desafíos}
    Demasiados estados, conocimiento difícil de codificar, reglas costosas de mantener y fragilidad ante excepciones.
  \end{block}
\end{frame}

\begin{frame}{Inviernos y resurgimiento estadístico}
  \begin{columns}[T]
    \column{0.48\textwidth}
      \begin{block}{Inviernos de la IA}
        \begin{itemize}
          \item Promesas desproporcionadas.
          \item Poca generalización.
          \item Cómputo y datos insuficientes.
          \item Sistemas frágiles.
        \end{itemize}
      \end{block}
    \column{0.48\textwidth}
      \begin{block}{Resurgimiento}
        \begin{itemize}
          \item Datasets extensos.
          \item Sensores y redes.
          \item Almacenamiento económico.
          \item Procesadores especializados.
        \end{itemize}
      \end{block}
  \end{columns}
\end{frame}

\begin{frame}{Deep learning e IA generativa}
  \begin{itemize}
    \item El aprendizaje profundo combina capas para aprender representaciones intermedias.
    \item Los modelos fundacionales se preentrenan sobre grandes colecciones.
    \item La IA generativa produce texto, imágenes, audio o código condicionados por una entrada.
  \end{itemize}
  \vspace{0.5em}
  \begin{alertblock}{Límite crítico}
    La fluidez es una propiedad de generación, no una prueba de verdad.
  \end{alertblock}
  Se requieren fuentes, validación, límites de uso y supervisión.
\end{frame}

\begin{frame}{El concepto de agente}
  Un agente:
  \begin{enumerate}
    \item recibe percepciones mediante sensores;
    \item mantiene o construye una representación;
    \item selecciona acciones;
    \item actúa mediante actuadores;
    \item recibe retroalimentación.
  \end{enumerate}
  Su función transforma una historia de percepciones en una acción:
  \[
    f: P^* \longrightarrow A
  \]
  El programa implementa esa función con memoria, conocimiento y cómputo limitados.
\end{frame}

\begin{frame}{Agente y entorno}
  \begin{center}
    \begin{tabular}{c c c}
      Entorno & $\xrightarrow{\quad\text{percepciones}\quad}$ & Agente \\
      Entorno & $\xleftarrow{\qquad\text{acciones}\qquad}$ & Agente
    \end{tabular}
  \end{center}
  \begin{block}{Ejemplo de movilidad}
    \begin{itemize}
      \item \concepto{Sensores}: GPS, reloj, demanda y tráfico.
      \item \concepto{Agente}: recomendador de reasignación.
      \item \concepto{Actuadores}: mensajes, órdenes o actualización de rutas.
      \item \concepto{Entorno}: flota, vías, pasajeros, clima y otros vehículos.
    \end{itemize}
  \end{block}
\end{frame}

\begin{frame}{PEAS: especificar la tarea}
  \small
  \begin{tabularx}{\textwidth}{>{\bfseries}p{0.25\textwidth} X}
    \toprule
    Componente & Pregunta \\
    \midrule
    Performance & ¿Cómo se mide el desempeño? \\
    Environment & ¿En qué entorno actúa? \\
    Actuators & ¿Cómo interviene en el entorno? \\
    Sensors & ¿Qué puede percibir? \\
    \bottomrule
  \end{tabularx}
  \vfill
  \begin{block}{Idea central}
    PEAS evita describir al agente solo por su algoritmo. Primero se especifican tarea y entorno; después se elige la arquitectura.
  \end{block}
\end{frame}

\begin{frame}{Propiedades del ambiente}
  Un ambiente puede ser:
  \begin{columns}[T]
    \column{0.48\textwidth}
      \begin{itemize}
        \item total o parcialmente observable;
        \item determinista o estocástico;
        \item episódico o secuencial;
        \item estático o dinámico;
      \end{itemize}
    \column{0.48\textwidth}
      \begin{itemize}
        \item discreto o continuo;
        \item de uno o múltiples agentes;
        \item conocido o desconocido.
      \end{itemize}
  \end{columns}
  \vfill
  \centering\alert{La clasificación depende del nivel de descripción y condiciona el diseño.}
\end{frame}

\begin{frame}{Ejemplo: recomendador de rutas}
  \scriptsize
  \begin{tabularx}{\textwidth}{>{\bfseries}p{0.20\textwidth} p{0.22\textwidth} X}
    \toprule
    Propiedad & Caracterización & Justificación \\
    \midrule
    Observabilidad & Parcial & No conoce todos los incidentes o intenciones \\
    Resultado & Estocástico & Una acción puede producir tiempos distintos \\
    Dependencia & Secuencial & La ruta actual modifica opciones posteriores \\
    Cambio & Dinámico & El tráfico cambia mientras se calcula \\
    Estado & Continuo & Tiempo, posición y velocidad varían \\
    Participantes & Multiagente & Interactúa con conductores y vehículos \\
    \bottomrule
  \end{tabularx}
\end{frame}

\begin{frame}{Caso transversal: movilidad urbana}
  \begin{block}{Necesidad}
    Reducir demoras mediante reasignación de vehículos entre zonas.
  \end{block}
  \begin{block}{Agente propuesto}
    Asistente operativo que observa demanda y flota, compara reasignaciones y recomienda una acción.
  \end{block}
  \begin{itemize}
    \item ¿Recomienda o ejecuta?
    \item ¿Cada cuánto vuelve a decidir?
    \item ¿Puede abstenerse?
    \item ¿Qué errores requieren intervención humana?
  \end{itemize}
\end{frame}

\begin{frame}{PEAS inicial para movilidad}
  \scriptsize
  \begin{tabularx}{\textwidth}{>{\bfseries}p{0.20\textwidth} X}
    \toprule
    Componente & Especificación inicial \\
    \midrule
    Desempeño & Demora, cancelaciones, distancia vacía, seguridad y equidad \\
    Entorno & Vías, zonas, flota, pasajeros, tráfico, clima y regulación \\
    Actuadores & Recomendar traslado, asignar vehículo, mantener o alertar \\
    Sensores & GPS, viajes, demanda, capacidad, tiempo y tráfico \\
    \bottomrule
  \end{tabularx}
  \vfill
  \pregunta{¿Qué ocurre si optimizamos solamente la demora promedio?}
\end{frame}

\begin{frame}{Actividad guiada: diseñar el agente}
  \small
  En equipos, completar:
  \begin{enumerate}
    \item propósito y usuario del agente;
    \item medida de desempeño con al menos tres criterios;
    \item entorno y actores afectados;
    \item acciones permitidas y acción de abstención;
    \item percepciones disponibles antes de decidir;
    \item propiedades del ambiente;
    \item riesgo de una métrica inadecuada;
    \item mecanismo de supervisión o falla segura.
  \end{enumerate}
  \begin{block}{Tiempo sugerido}
    30 minutos de trabajo + 10 minutos de contraste.
  \end{block}
\end{frame}

\begin{frame}{Criterios de revisión}
  La especificación es adecuada si:
  \begin{itemize}
    \item las percepciones existen antes de elegir la acción;
    \item los actuadores producen cambios o recomendaciones claras;
    \item el desempeño combina objetivo, costos y restricciones;
    \item la caracterización del entorno está justificada;
    \item se distingue recomendación de ejecución autónoma;
    \item existe una respuesta segura ante incertidumbre o fallos;
    \item se identifican responsables y personas afectadas.
  \end{itemize}
  \begin{block}{Entregable de semana 1}
    Sección PEAS de la ficha del proyecto.
  \end{block}
\end{frame}

\begin{frame}{Puesta en común}
  Cada equipo presenta en un minuto:
  \begin{enumerate}
    \item la acción principal del agente;
    \item su medida de desempeño;
    \item una propiedad crítica del ambiente;
    \item una condición en la que debe abstenerse.
  \end{enumerate}
  \begin{block}{Preguntas para contrastar}
    \begin{itemize}
      \item ¿Lo racional según la métrica es aceptable para los usuarios?
      \item ¿Qué percepción faltante genera más incertidumbre?
      \item ¿Qué cambia si el agente ejecuta en vez de recomendar?
    \end{itemize}
  \end{block}
\end{frame}

\begin{frame}{Síntesis}
  \begin{itemize}
    \item ``Inteligencia'' puede evaluarse como conducta, cognición, inferencia o acción racional.
    \item El test de Turing evalúa comportamiento, no verdad ni inteligencia general.
    \item La historia relaciona representaciones, datos, cómputo y expectativas.
    \item Un agente transforma percepciones en acciones dentro de un entorno.
    \item La racionalidad depende de información, objetivos y restricciones.
    \item PEAS permite especificar una tarea antes de elegir el algoritmo.
  \end{itemize}
\end{frame}

\begin{frame}{Autoevaluación}
  \small
  \begin{enumerate}
    \item ¿Qué diferencia existe entre actuar como humano y actuar racionalmente?
    \item ¿Por qué pensar racionalmente no garantiza actuar correctamente?
    \item ¿Qué limitación principal tiene el test de Turing?
    \item ¿Qué favoreció el resurgimiento estadístico de la IA?
    \item ¿Qué componentes forman una especificación PEAS?
    \item ¿Por qué el tráfico es parcialmente observable y dinámico?
    \item ¿Cómo puede un agente ejecutar correctamente una función mal definida?
  \end{enumerate}
\end{frame}

\begin{frame}{Lecturas y recursos}
  \begin{block}{Lectura principal}
    \begin{itemize}
      \item \textbf{Capítulo 1}: secciones 1.1.2, 1.2 y 1.3.
      \item \textbf{Capítulo 6}: secciones 6.1 y 6.2.
    \end{itemize}
  \end{block}
  \begin{block}{Para continuar}
    \begin{itemize}
      \item Revisar el glosario y la autoevaluación del capítulo 1.
      \item Integrar PEAS con la ficha elaborada en Data Science.
    \end{itemize}
  \end{block}
\end{frame}

\end{document}
